% =====================================================================
%  c-intro.tex —— 原文 pp.4--7（Lead）
%  内容：Introduction and background Bibliography（引言与背景·参考书目）
%  原文 pp.1--3 为标题页与目录，不逐字翻译（见 c00.tex）。
% =====================================================================
% ===== 原文 pp.4--7 =====
% [原文 p.4]
\section*{引言与背景}
\addcontentsline{toc}{section}{引言与背景}

1995 年，伦敦数学会（London Mathematical Society, LMS）在兰卡斯特大学（University of Lancaster，
英国）组织了一期 K 理论研究生暑期学校。这期暑期学校覆盖了 K 理论的三个基本分支，即
拓扑 K 理论、解析 K 理论（$K$-同调）与高阶代数 K 理论。主讲人分别是
J. D. S. Jones 教授（Warwick）、J. Roe 教授（当时在 Oxford，现在 Penn. State）与
D. G. Quillen 教授（Oxford）。

本书作者当时是牛津大学数学研究所的研究生，有幸成为该暑期学校的一名学员，他的个人笔记构成了本书的骨架。
此后数年间，作者两度讲授 K 理论研究生课程（分别在英国的 Cardiff 大学与希腊的 Athens 大学），
这些笔记也随之扩充与打磨。这一经历促使作者把全书组织成 18 个小时报告的形式，每个分支 6 讲。
后来他又补写了一章，介绍 Waldhausen K 理论（不过这一章从未在课堂上讲过），
另加两节：一节讲扭曲 K 理论及其在弦论/M-理论中的应用，一节讲固态物理中所谓能隙标记（gap labeling）问题。
为方便读者，我们还附了一个附录，收录数学其他领域（尤其是代数拓扑）的若干结论。
有少数命题没有给出证明，这或是因为证明相当直接（这种情况留给读者作练习），
或是因为证明冗长复杂，我们只展示主要想法并给出参考文献。

% [原文 p.5]
就 K 理论而言，数学文献不能算丰富。我们给出了一份几乎完整的相关书目（本书的笔记对这些书有所取用）：

拓扑 K 理论方面，有 M. F. Atiyah 1967 年经典的哈佛讲义（见 M.F. Atiyah: "K-Theory", Benjamin 1967），
以及 Efton Park 的一本新书："Complex Topological K-Theory", Cambridge University Press 2008。
此外还有康奈尔大学 A. Hatcher 的一些网上讲义（见 http://www.math.cornell.edu/ hatcher/），
以及巴黎 Max Karoubi 的一些网上讲义（见 http://people.math.jussieu.fr/ karoubi/KBook.html），
连同他的著作 "K-Theory - An Introduction", Springer (1978)。

解析 K 理论方面，有 N. E. Wegge-Olsen 的书："K-Theory and C$^*$ Algebras", Oxford University Press 1993；
B. Blackadar 更进阶级别的书："K-Theory for Operator Algebras", Cambridge University Press 1998；
Nigel Higson 与 John Roe 的书：Analytic K-Homology, Oxford University Press 2000；
以及 F. Larsen、M. Rørdam 与 M. Rørdam 的书："An Introduction to K-Theory for C$^*$-algebras",
Cambridge University Press 2000。还有剑桥大学 T. Gowers 的一些不错的简短网上讲义：
"K-Theory of Banach Algebras"（见 http://www.dpmms.cam.ac.uk/ wtg10/）。

高阶代数 K 理论方面，有 Jonathan Rosenberg 的书："Algebraic K-Theory and its Applications",
Springer Graduate Texts in Mathematics (1994)；V. Srinivas 的书："Algebraic K-Theory", Birkhauser, Boston (1996)；
以及 Hvedri Inassaridze 的书："Algebraic K-Theory", Kluwer 1995。
此外还有 C. Weibel 的一些网上讲义："The K-book: An introduction to Algebraic K-Theory"。

由此可见，文献中并没有一本同时覆盖 K 理论三个基本分支的简明导引。

K 理论之父是 M. F. Atiyah 与 A. Grothendieck，两人可以说是 20 世纪下半叶最伟大的数学家。
这个名称是 Grothendieck 起的，字母 "$K$" 代表德语词 "(die) Klasse"，英语意为"类"。
K 理论是所谓广义（或"异域"）同调理论之一，满足 Eilenberg--Steenrod 同调公理五条中的四条：
同伦、切除、可加与正合，但不满足维数公理——点的 K 理论是 $\Z$ 而不是零。
在其诞生以来的 50 年间，K 理论已被证明是许多数学领域的有用工具，包括拓扑、整体分析、
指标理论、数论、代数、（非交换）几何等。此外 K 理论在理论物理中也找到了许多应用，
例如量子场论中的反常抵消、固态物理中所谓能隙标记问题、M-理论中膜的拓扑荷
（其中膜的拓扑荷被视为扭曲 K 理论的类）等等。

% [原文 p.6]
K 理论在数学上的演进包括 M. F. Atiyah 与 G. Segal 的等变 K 理论、用于流形手术的 $L$ 理论、
$KK$ 理论（Kasparov K 理论或双变 K 理论）、A. Connes 的 $E$ 理论、
Waldhausen K 理论或"$A$ 理论"（它是 Quillen 高阶代数 K 理论的拓扑版本）等。
我们还应提到高阶代数 K 理论与 V. Voevodsky 的动机（motivic）理论（动机上同调）之间的密切关系。

本书适合研究生阅读，因此我们假定读者具备良好的代数知识（尤其是群、模与同调代数）、
几何知识（流形理论）与代数拓扑知识（CW 复形的奇异（上）同调与链（上）同调，以及基本同伦论）。
证明的结束用一个空心方块（$\square$）表示。在记号上，例如命题 2.3.9 指第 2 章第 3 节（讲）中的第 9 个命题。
下面给出一些合适的背景推荐书目：

\para{代数}
\begin{itemize}[leftmargin=1.6em,itemsep=1pt]
  \item S. MacLane, G. Birckhoff: "Algebra", Chelsea, 1988
  \item S. Lang: "Algebra", Addison Wesley 1993
  \item P. J. Hilton, U. Stammbach: "A Course in Homological Algebra", Springer 1997.
  \item N. Bourbaki, "Algebra", Vol I, II
  \item W. A. Adkins, S. H. Weintraub, "Algebra", Springer, 1992
  \item P. M. Cohn "Algebra", John Wiley, 1989
  \item T. H. Hungerford, "Algebra", Springer, 1980
\end{itemize}

\para{几何}
\begin{itemize}[leftmargin=1.6em,itemsep=1pt]
  \item N. Hitchin 关于流形的网上讲义，见
    http://people.maths.ox.ac.uk/ hitchin/hitchinnotes/hitchinnotes.html
  \item P. Griffiths and J. Harris: "Principles of Algebraic Geometry", Wiley, 1994
  \item S. Kobayashi, K. Nomizu: "Foundations of Differential Geometry", Vol I, II, Wiley, 1996
  \item J. M. Lee, "Introduction to Smooth Manifolds", Springer, 2002
\end{itemize}

% [原文 p.7]
\para{拓扑}
\begin{itemize}[leftmargin=1.6em,itemsep=1pt]
  \item A. Hatcher，代数拓扑网上讲义，见 http://www.math.cornell.edu/ hatcher/AT/AT.pdf
  \item J. P. May: "A concise Course in Algebraic Topology"，
    见 http://www.math.uchicago.edu/ may/CONCISE/ConciseRevised.pdf
  \item E. H. Spanier: "Algebraic Topology", Springer 1991
  \item A. Dold: "Lectures on Algebraic Topology", Springer 1991
  \item C. R. F. Maunder: "Algebraic Topology", Dover 1996
  \item W. S. Massey: "A Basic Course in Algebraic Topology", Springer 1991
  \item J. Munkres: "Elements of Algebraic Topology", Westview 1995
  \item R. Bott, L. Tu: "Differential forms in algebraic topology", Springer 1982
  \item N. Steenrod: "The topology of fibre bundles", Princeton (1951)
  \item S. Eilenberg and N. Steenrod: "Foundations of Algebraic Topology", Princeton (1952)
\end{itemize}

\noindent IPZ，Athens，2010 年 8 月\\
http://sites.google.com/site/ipzoisscience/home

% [Tutor Note: p.6 原文连字符处印作 "Birckoff"，是 Birckhoff（MacLane--Birkhoff 的《Algebra》）的排印讹误，正文照原样排。
% p.5 原文 "Eilenberg-Steenroad" 应为 Eilenberg--Steenrod（同页下方另有一处写作 Steenrod），正文照原样排。
% p.5 原文 "Cardiff" 为 Cardiff（连字 ff 丢失所致，印刷本身无误）。
% p.5 "Rrdam" 实为 Rørdam（同句前一作者已写作 Rørdam），原文排印漏了 ø，正文照原样排。
% p.4 "the the Grothendieck Group"（p.3 目录）与正文多处 "clasic""setions""cercises" 等拼写错误，均照原文排，
% 不在此处逐一登记，只在此说明：本译本一律不"顺手改正"原文拼写。
% 另外：原文 p.5 的 URL 中 "hatcher/"、"karoubi/"、"wtg10/" 前应有一个波浪号（~），PDF 字形层丢失，译本按原样保留空格。]
